\documentclass[10pt,a4paper]{amsart}
\usepackage[margin=19mm]{geometry}
\usepackage{amsmath,amssymb,mathtools}
\usepackage[scr=rsfs,cal=euler]{mathalfa}
\usepackage{xfrac}
\usepackage[unicode,hidelinks,bookmarks=false]{hyperref}
\newcommand{\ie}{i.e.\ }
\newcommand{\cf}{cf.\ }
\newcommand{\resp}{respectively\ }
\newcommand{\Subsection}{\subsection}
\providecommand{\nobreakdash}{\nobreak\hbox{-}}
\setlength{\emergencystretch}{2em}
\multlinegap0pt
\usepackage{mathtools}
\usepackage{amssymb}
\usepackage{algorithm}\usepackage{algpseudocode}
\newtheorem{lem}{Lemma}
\title{Exponential stability of the {linearized} viscous Saint-Venant equations using a quadratic Lyapunov function}
\author{Amaury Hayat, Nathan Lichtl\'e}
\date{Source: arXiv:2603.06411v1 (CC BY 4.0)}
\begin{document}
\maketitle

\noindent\emph{Source and licence.} Exponential stability of the {linearized} viscous Saint-Venant equations using a quadratic Lyapunov function. Version arXiv:2603.06411v1 (CC BY 4.0), distributed by arXiv under the Creative Commons Attribution 4.0 International licence, http://creativecommons.org/licenses/by/4.0/, as recorded in the arXiv licence metadata for this exact version and shown on the abstract page https://arxiv.org/abs/2603.06411v1. Full licence notice: \texttt{licenses/arxiv-2603.06411-CC-BY-4.0.txt}.\par\smallskip

\noindent\textit{Editorial scope.} Counted unit: the article's lem lem:appendix (source0.tex lines 877--894). The article's complete original proof of it is delivered with it (source0.tex lines 895--946, one \texttt{proof} environment of 52 lines). No mathematics is written, repaired or re-derived: the statement and the proof are the article's own lines, in the article's own notation, and no retained line is substituted or re-flowed.  Retained uncounted as the source-local context the proof needs: lines 851--854; lines 869--875.  Nothing was omitted from any retained span, so the package carries no omission record.  No retained line contains a citation, so the wrapper carries no bibliography and no bibliography entry.  No retained line contains a figure environment, an image-inclusion command or drawing code, so no image is delivered and the package declares no asset dependency.  Editorial formatting only, in addition: an AMS \texttt{amsart} preamble that restates the article's own declarations and macros used by the retained lines, namely the environments \texttt{lem} on separate counters, together with the standard packages those lines need.  The retained environments print in this document as Lemma 1 in order of appearance, whereas the article prints the same items with the numbering of its own full document.  The complete native source of the article as distributed in the arXiv e-print for this exact version is bundled unchanged as \texttt{sources/195867/source0.tex}.
\begin{equation}
\label{eq:deftildef}
\tilde{f}(V^{*}) = \frac{{V^{*}}^{2}\kappa}{3\mu V^{*}+\kappa Q^{*}}.
\end{equation}

\begin{equation}
\label{eq:sys4}
    \begin{cases}
        &y_{x} = z\\
        & \mu z_{x} = \frac{1}{4}z\left(y-g \frac{Q_{0}}{y^{2}}\right)+\frac{3}{4}\mu \frac{\tilde{f}(y)y}{Q_{0}}+\mu \frac{z^{2}}{y}.
    \end{cases}
\end{equation}

\begin{lem}
\label{lem:appendix}
Let $Q_{0}>0$, $\varepsilon>0$, $c_{y}>0$, $C_{z}>0$ and
\begin{equation}
\begin{split}
\Omega(\varepsilon) := \Big\{&(y,z)\in (0,+\infty)\times\mathbb{R} \quad \Big| \quad \frac{gQ_{0}}{y^{2}}-y>\varepsilon,\; y\geq c_{y},\; |z|\leq C_{z} \Big\}.
\end{split}
\end{equation}
There exists $\mu^{*}>0$ such that for any $\mu\in(0,\mu^{*})$ and any $y_{0},z_{0}\in \Omega(2\varepsilon)$, there exists a unique solution $(y,z)\in C^{1}([0,L];\Omega(\varepsilon))$ to \eqref{eq:sys4} with initial condition $(y_{0},z_{0})$. Moreover there exists $C_{1}>0$ depending only on $\varepsilon$, $Q_{0}$, $c_{y}$ and the parameters of the system such that 
\begin{equation}
\label{eq:estimyz}
\begin{split}
|y(x)-y_{0}|&\leq \frac{4\mu}{\varepsilon}\left(1-e^{-\frac{\varepsilon x}{4\mu}}\right)|z_{0}| + C_{1}\frac{4\mu}{\varepsilon}\left(x-\frac{4\mu}{\varepsilon}(1-e^{-\frac{\varepsilon x}{4\mu}}))\right),\\
|z(x)|&\leq  |z_{0}| e^{-\frac{\varepsilon x}{4\mu}}+ C_{1}\frac{4\mu}{\varepsilon}(1-e^{-\frac{\varepsilon x}{4\mu}}),
\end{split}
\end{equation}
{and for any $x_{0}\in[0,L]$ such that $z(x_{0})=0$, $z_{x}(x_{0})>0$.}
\end{lem}

\begin{proof}
Let $Q_{0}>0$, $\varepsilon>0$, $c_{y}>0$, $C_{z}>0$ and let $\mu\in(0,\mu^{*})$ with $\mu^{*}$ to be chosen later on. {Note that on $\Omega(\varepsilon)$, the function $(y,z)\mapsto (z, z(y-gQ_{0}/y^{2})/4+3\mu\tilde{f}(y)y/(4Q_{0})+\mu z^{2}/y)$ is of class $C^{1}$ and hence locally Lipschitz. Thus from} the {(local)} Cauchy-Lipschitz theorem, 
there exists a unique maximal solution $(y, z)$ in $\Omega(\varepsilon)$ to \eqref{eq:sys4} defined 
on $[0,L^{*})$ with $L^{*}=+\infty$ or $L^{*}\in(0,+\infty)$ and in {the second} case either
\begin{equation}
\label{eq:boundslim}
\begin{split}
&\lim\limits_{x\rightarrow L^{*}}(gQ_{0}/y^{2}(x)-y(x)) =\varepsilon\text{ with }y(x)> 0\\
\text{ or }&\lim\limits_{x\rightarrow L^{*}}y(x) = c_{y}\\
\text{ or }&\lim\limits_{x\rightarrow L^{*}}|z(x)| = C_{z}
\end{split}
\end{equation}
Note that $gQ_{0}-y^{3}-\varepsilon y^{2}$ has a unique non-negative (and actually positive) root (because it is positive at $y=0$, negative when $y \to +\infty$ and its derivative is negative for $y > 0$) that we denote by $y_{1,\varepsilon}>0$, this means that 
\begin{equation}
\label{eq:defy1}
\lim\limits_{x\rightarrow L^{*}}(gQ_{0}/y^{2}(x)-y(x)) =\varepsilon \iff
   \lim\limits_{x\rightarrow L^{*}}y(x) = y_{1,\varepsilon}.
\end{equation}
In any cases, on $[0,L^{*})$, we can define 
\begin{equation}
\label{eq:defg}
{\Gamma(x)} =\frac{gQ_{0}}{y^{2}(x)}-y(x),
\end{equation}
and we have
\begin{equation}
    \mu z'(x) = -\frac{1}{4}{\Gamma(x)}z(x)+\mu\left(\frac{3}{4}\frac{\tilde{f}(y(x))y(x)}{Q_{0}}+\frac{z(x)^{2}}{y(x)}\right),
\end{equation}
hence, by Duhamel's formula,
\begin{equation}
\begin{split}
    z(x) &= z_{0} e^{-\frac{1}{4\mu}\int_{0}^{x}{\Gamma(s)}ds}+  \int_{0}^{x}e^{-\frac{1}{4\mu}\int_{s}^{x}{\Gamma(\tau)}d\tau}\left(\frac{3}{4}\frac{\tilde{f}(y(s))y(s)}{Q_{0}}+\frac{z(s)^{2}}{y(s)}\right) ds.
\end{split}
\end{equation}
Since $(y(s),z(s))\in \Omega(\varepsilon)$ for any $s\in [0,x]\subset[0,L^{*})$ and from \eqref{eq:defy1},
\begin{equation}
\label{eq:estimz}
\begin{split}
    |z(x)| &\leq |z_{0}| e^{-\frac{\varepsilon x}{4\mu}}+C_{1}\int_{0}^{x}e^{-\frac{\varepsilon (x-s)}{4\mu}} ds \\&= |z_{0}| e^{-\frac{\varepsilon x}{4\mu}}+ C_{1}\frac{4\mu}{\varepsilon}(1-e^{-\frac{\varepsilon x}{4\mu}}),
\end{split}
\end{equation}
where $C_{1} = 3Q_{0}^{-1}\max\limits_{[c_{y},y_{1,\varepsilon}]}|\tilde{f}(y)y|/4+ C_{z}^{2}c_{y}^{-1}$. 
Therefore, using the first equation of \eqref{eq:sys4},
\begin{equation}
\label{eq:estimy}
\begin{split}
|y(x) - y_{0}| &\leq \int_{0}^{x} |z_{0}| e^{-\frac{\varepsilon s}{4\mu}}+ C_{1}\frac{4\mu}{\varepsilon}(1-e^{-\frac{\varepsilon s}{4\mu}}) ds\\
&= \frac{4\mu}{\varepsilon}\left(1-e^{-\frac{\varepsilon x}{4\mu}}\right)|z_{0}|  + C_{1}\frac{4\mu}{\varepsilon}\left(x-\frac{4\mu}{\varepsilon}(1-e^{-\frac{\varepsilon x}{4\mu}}))\right).
\end{split}
\end{equation}
We see from \eqref{eq:estimz}--\eqref{eq:estimy} that we can choose $\mu^{*}$, depending on $C_{z}$, $c_{y}$ and $\varepsilon$ such that $L^{*}>L$ (otherwise we get a contradiction because we never reach \eqref{eq:boundslim}). {In fact, this characterizes $\mu^{*}$ and it can be computed as the limiting constant at which \eqref{eq:boundslim} is reached.} 
This shows the existence (and unicity) of the solution $(y,z)\in C^{1}([0,L];\Omega(\varepsilon))$. In addition the estimates \eqref{eq:estimz}--\eqref{eq:estimy} hold, {which is exactly \eqref{eq:estimyz}. Finally, assume that for $x_{0}\in[0,L]$, $z(x_{0})=0$, since $(y(x_{0}),z(x_{0}))\in \Omega(\varepsilon)$, $y(x_{0})\geq c_{y}>0$ and, from \eqref{eq:sys4} and the expression of $\tilde{f}$ given by \eqref{eq:deftildef} we deduce directly that $z_{x}(x_{0})>0$,} which ends the proof of Lemma \ref{lem:appendix}.
\end{proof}

\end{document}
